\section{The bridge, and the wall}
\label{sec:bridge}

This is the one section that addresses the motivating question of
Section~\ref{sec:intro}. We hold to the firewall: the mathematics below stands on its own;
the physical reading is labelled and cites the mathematics one way. The result is a pair
--- a forced symmetry bridge and a confirmed dimensional wall --- and we tag each
component as \sForced\ (a mathematical identity), \emph{permitted} (consistent but not
forced), or \emph{rhyme} (a suggestive homonym that is \emph{not} the bridge).

\subsection{The bridge is real}

The once-punctured-torus theory in the class-$S$ construction is the
$\mathcal{N}=2^\ast$ theory (maximally supersymmetric Yang--Mills deformed by an adjoint
mass). Its Coulomb-branch geometry is governed by the \emph{same} Fricke cubic
\eqref{eq:fricke}: the character variety of the once-punctured torus, with the very trace
coordinates of Section~\ref{sec:invariant}. Crucially, the $S$-duality group acts on that
character variety \emph{through the mapping class group}, not merely on the ultraviolet
gauge coupling $\tau$ \cite{allegrettishan2024}.

\begin{theorem}[\sForced, literature-confirmed]\label{thm:bridge}
The $\SL(2,\Z)$ trace-map action of the metallic monodromies on the Fricke character
variety (Proposition~\ref{prop:invariant}) coincides with the $S$-duality /
mapping-class action of the $\mathcal{N}=2^\ast$ class-$S$ theory on its Coulomb branch.
The metallic conjugacy classes, the dilatation eigenvalue $\lambda_m^2$ (the Cantat--Loray
dynamical degree \cite{cantatloray2009}), the fixed slopes, and the Markov fibre are all
reproduced, because it is the same group acting the same way on the same variety
\cite{allegrettishan2024,gmn2013}.
\end{theorem}

This is the widest reach the one-object picture achieves, and it is \sForced, not a
rhyme. The honest caveats are equally precise, and we name them in the same vocabulary:
\begin{itemize}[leftmargin=1.4em]
\item \emph{Forced:} the action on the \emph{character variety} reproduces the anchors of
  Section~\ref{sec:invariant}. This is the bridge.
\item \emph{Rhyme:} ``$\SL(2,\Z)$ acts modularly on the gauge coupling $\tau$'' is the
  well-known $S$-duality \emph{homonym}; it acts on a different space and is \emph{not}
  the trace-map dynamics. We flag it as rhyme precisely so it is not mistaken for the
  bridge.
\item \emph{Rhyme / permitted:} the gauge group is free input and $\mathcal{N}=2^\ast$ is
  non-chiral, so at the level of physical gauge content there is no forced identification.
\end{itemize}
Even the forced bridge is mathematics: it is a symmetry identity, not a crossing to
physics.

\subsection{The wall is confirmed}

That the bridge is real is exactly why the firewall had to be tested hardest: a genuine
symmetry bridge makes a scale-crossing \emph{feel} plausible. So we asked the decisive
question. In the 3d--3d correspondence applied to the figure-eight, does any quantity
carry a \emph{scale} across --- does the complex volume / Chern--Simons invariant enter
the partition function as anything other than a dimensionless exponent?

\begin{theorem}[\sKnown, literature reading]\label{thm:wall}
The complex volume $\widehat c(\rho)=i(\Vol+i\,\CS)$ of an $\SL(n,\C)$ representation is a
dimensionless element of $\C/4\pi^2\Z$ \cite{gtz2015}; in the 3d--3d state integral it
enters only as the exponent saddle, with all dimensionful content carried by the level
$\hbar\leftrightarrow k$ and the squashing radius \cite{dimofte2014,cordovajafferis2013}.
For the figure-eight, $\Vol=2.029883\ldots$ and $\CS=0$: the candidate scale-carrier
vanishes identically. No $\kk$-type or volume-type invariant of the object can source a
physical scale (\fig{fig:bridge}).
\end{theorem}

\begin{figure}[t]
  \centering
  \includegraphics[width=\linewidth]{fig12_bridge_and_wall}
  \caption{The bridge and the wall. \emph{Left (forced):} the metallic $\SL(2,\Z)$
  trace-map symmetry \emph{is} the $\mathcal{N}=2^\ast$ class-$S$ $S$-duality /
  mapping-class action on the Fricke character variety. \emph{Right (confirmed):} the
  complex volume is dimensionless ($\C/4\pi^2\Z$); the scale lives in $\hbar/k$ and the
  squashing radius, and for the figure-eight the scale-carrier $\CS$ vanishes. A symmetry
  identity is not a scale. \sForced\,/\,\sKnown}
  \label{fig:bridge}
\end{figure}

The blade, held to the end: a \emph{symmetry} identification is not a \emph{scale}. ``The
partition function has units'' is not a crossing --- the units are in $\hbar/k$ and the
radius. A crossing would require a dimensionful quantity to attach to the \emph{invariant
itself}, and the literature exhibits the opposite.

\begin{remark}[\sPost]\label{rem:postulated}
There is a sharper, \emph{structural} reason one might give for the wall: a conserved /
topological invariant cannot run, and running is what manufactures a scale; so an object
whose content is a topological invariant cannot, in principle, source one. This is
compelling and was proposed in review, but it has not been put through the project's
verification gate, and so we record it as \sPost\ --- not as the banked result. The banked
wall is the literature reading of Theorem~\ref{thm:wall}.
\end{remark}

\subsection{The seam: the object is open, and closing it is where structure actualizes}

The wall above is about \emph{scale}. A second reading sharpens \emph{where the object's
symmetries break} --- and it turns on a fact we had under-used: the figure-eight is an
\emph{open} object. It is a complement $S^3\setminus K$, defined by what is removed, with a
cusp (boundary torus) as its interface with the exterior. Its amphichirality and
CP-symmetry are \emph{open}-object symmetries; \emph{closing} it --- a Dehn filling --- is
where they break. (Curie's principle, sometimes invoked to argue a symmetric object cannot
self-select an asymmetric outcome, is a \emph{closed}-system statement; it does not bind a
complement, which breaks symmetry through boundary conditions.)

\begin{proposition}[\sKnown\ mathematics, \sForced\ identification]\label{prop:seam}
The figure-eight is the once-punctured-torus bundle with monodromy $A=\big(\begin{smallmatrix}2&1\\1&1\end{smallmatrix}\big)$.
Among its ten exceptional Dehn fillings, the fiber ($0$-slope) filling is the \emph{unique}
$\mathrm{Sol}$ torus bundle, with monodromy exactly $A$ --- the trace-map matrix $A=LR$ of
\S\ref{sec:object} and its torsion ladder $\lvert\det(A^n-I)\rvert$. Every generic filling
is chiral ($\CS\notin\{0,\tfrac12\}$), with the sign law $\CS(p,-q)=-\CS(p,q)$; the
geometric handedness realises complex conjugation of the $\Q(\sqrt{-3})$ trace field --- the
same Galois involution that flips the commutator phase $\kappa=u^2+2=\sqrt3\,e^{\mp i\pi/6}$
between $\pm\pi/6$. The core geodesic of the $(1,n)$ filling has length $\sim 2\pi/n$ (a
Neumann--Zagier ladder); the peripheral $(\mu,\lambda)$ pairing is the Goldman/NZ symplectic
form, a filling a choice of polarization.
\end{proposition}

The mathematics here is classical (Thurston's exceptional surgeries; Garoufalidis--Jeon on
trace-field growth; Goldman and Neumann--Zagier on the symplectic structure); the new part
is the \emph{reading}, and it cuts in a definite direction. Tested across the fillings, the
seam is \emph{selective for the object's own structure} --- the canonical closing re-sees
$A=LR$ --- but a \emph{catalogue for Standard-Model values}: the $E_6$/$\Q(\sqrt{-3})$
arithmetic is \emph{lost} on closing (no filling re-instantiates it, sharpening the
input-$E_6=$output-$E_6$ question toward ``catalogue''), the CP \emph{sign} stays external
(the object is CP-symmetric: it delivers $\pm\pi/6$ symmetrically), the scale \emph{value}
is gapped, and the chiral and dynamical data are not supplied.

\begin{remark}[\sPost]\label{rem:structural-theorem}
This relocates the project's structural theorem to the seam: \emph{the object supplies a
canonical closing and all the dimensionless structure --- the dynamics $A=LR$, the CP
magnitude $\pi/6$ with its forced sign law, the scale ladder $2\pi/n$, the
peripheral-symplectic clock --- while the actualization of a specific Standard-Model-valued
world is not forced}. It is the architecture-not-furniture stance made mechanistic at the
boundary, and it \emph{strengthens} the wall rather than crossing it. As a physics reading
it is \sPost; the mathematics it cites is banked and firewalled.
\end{remark}

\subsection{Where it reaches, and where it stops}

The physics arc closes in the strongest honest place available: a real bridge --- the
object's symmetry \emph{is} the duality action of a gauge theory, forced and
literature-confirmed --- and a confirmed wall --- that bridge carries structure, not
scale. The cosmological-constant question lies on the far side of a wall this structure
does not cross. That pairing is not a failure; it is a result. The work reached across the
widest gap it could, mapped exactly where it arrives and where it stops, and exhibited the
reasons rather than asserting them. What remains after this is mathematics, not physics ---
the subject of the next section.
